\section{ARC Construction from Blind Signatures}
\label{sec:arc-construction}

This section turns the blind signature of Section~\ref{sec:blind-signature} into an anonymous rate-limited credential
(ARC) system. We use the following terminology throughout:
\emph{issuance} refers to the blind signature protocol, while \emph{showing} refers to an ARC presentation.
The issuance proof is the \emph{pre-sign NIZK}; the showing proof is the \emph{post-sign NIZK}.

\subsection{From blind signatures to ARC}
A credential issued to a holder conceptually consists of:
\begin{itemize}[leftmargin=1.25em]
  \item the signed message \(m=\mu\Vert \veck\), where \(\mu\) are attributes and \(\veck\) is a uniformly sampled secret PRF key;
  \item the signature \(u\) satisfying \(A u = h_{m,(r_0,r_1)}(B)\) for some bounded randomness \((r_0,r_1)\).
\end{itemize}
The holder may additionally store local auxiliary values (e.g., the randomness \((r_0,r_1)\)) needed to prove the
signature relation in zero knowledge. These values are never revealed in showings.

\paragraph{Message structure.}
Every credential message has the abstract form
\[
  m \;=\; \mu \Vert \veck.
\]
Here \(\mu\) encodes the holder attributes and \(\veck\) is sampled uniformly and serves as the secret PRF key used for
rate limiting. In the implementation this abstract message is stored in packed witness rows \((M,K)\), but that is
only a representation choice and not part of the abstract syntax.

\subsection{Rate-limiting via the secret key \(\veck\)}
Rate limiting is implemented by deterministic tags derived from a PRF keyed by the credential’s secret \(\veck\). For a
public nonce \(\nonce\in\mathcal{D}\), the holder computes
\[
  \prftag \;=\; \PRF(\veck,\nonce).
\]
Intuitively:
\begin{itemize}[leftmargin=1.25em]
  \item If the holder never repeats a nonce, then \(\prftag\) is pseudorandom and yields unlinkable showings.
  \item If the holder repeats a nonce, the tag repeats deterministically, enabling the verifier to detect
  double-spending for that nonce.
\end{itemize}
The verifier enforces its rate policy by restricting acceptable nonces (e.g., encoding epoch, service, or a bounded
counter) and maintaining state that records previously accepted \((\nonce,\prftag)\) pairs for a given
policy context.

\subsection{Concrete PRF instantiation (Poseidon2-like)}
\label{subsec:poseidon-prf}

\providecommand{\Tr}{\mathrm{Tr}}

We instantiate the PRF from Definition~\ref{def:prf} using the Poseidon2-like construction of
Definitions~\ref{def:poseidon-prf} and \ref{def:poseidon-perm}. Algorithm~\ref{alg:prf} summarizes the computation.

\begin{algorithm}[H]
\caption{$\PRF(k,n)$ (Poseidon2-like with feed-forward + truncation)}
\label{alg:prf}
\begin{algorithmic}[1]
\State $x \gets k\Vert n$; $x^{(0)}\gets x$
\For{$i=1$ to $R_F/2$} \State $x\gets E_i(x)$ \EndFor
\For{$i=1$ to $R_P$} \State $x\gets I_i(x)$ \EndFor
\For{$i=R_F/2+1$ to $R_F$} \State $x\gets E_i(x)$ \EndFor
\State $y \gets x + x^{(0)}$ \Comment{feed-forward}
\State \Return $\Tr(y)$ \Comment{first $\ell_{\prftag}$ lanes}
\end{algorithmic}
\end{algorithm}

\paragraph{Proof-friendly arithmetization.}
In the post-sign NIZK we enforce $\prftag=\PRF(\veck,\nonce)$ using a checkpointed arithmetization: we commit
selected S-box outputs and let the verifier replay all linear wiring between checkpoints from openings and public
constants. Section~\ref{subsec:prf-checkpointed} states the constraint families; Appendix~\ref{app:prf} records the
checkpoint schedule and extended notes.

\subsection{Showing protocol}
\label{subsec:showing-protocol}

\begin{algorithm}[H]
\caption{\(\Show\): ARC presentation (canonical protocol)}
\label{alg:showing}
\begin{algorithmic}[1]
\Require Holder has credential \((\mu,\veck,r_0,r_1,u)\). Public parameters include \(A,B,\BoundMsg\) and PRF description.
\Require Input nonce \(\nonce\in\mathcal{D}\) that has not been used before for this credential.
\Ensure Presentation \((\prftag,\nonce,\pi_{\mathsf{post}})\).
\State Compute \(\prftag\gets \PRF(\veck,\nonce)\).
\State Produce a \emph{post-sign} NIZK \(\pi_{\mathsf{post}}\) proving knowledge of \((u,\mu,\veck,r_0,r_1)\) such that:
\begin{enumerate}[leftmargin=1.6em]
  \item \(A u = h_{\mu\Vert \veck,(r_0,r_1)}(B)\);
  \item \(\prftag = \PRF(\veck,\nonce)\);
  \item \(\mu,\veck,r_0,r_1\) are coefficient-wise in \([-\BoundMsg,\BoundMsg]\) and \(u\) satisfies the public shortness bound.
\end{enumerate}
\State Output \((\prftag,\nonce,\pi_{\mathsf{post}})\).
\end{algorithmic}
\end{algorithm}

The verifier runs \(\Verify\) by checking \(\pi_{\mathsf{post}}\) and consulting a policy state (e.g., a database) that
tracks used nonces/tags. The verifier rejects if the proof fails or if the nonce/tag violates the policy (e.g., already
spent, out of range, or malformed).

\subsection{PRF role and rationale}
The PRF serves two simultaneous roles:
\begin{itemize}[leftmargin=1.25em]
  \item \textbf{Public rate-limiting handle.} The tag is deterministic in \((\veck,\nonce)\) so the verifier can
  enforce stateful policies without interacting with the issuer and without learning the credential message.
  \item \textbf{Unlinkability across nonces.} For fresh nonces, tags are pseudorandom even given previous tags, so a
  verifier cannot link two presentations to the same credential unless the policy context forces reuse.
\end{itemize}
Crucially, \(\veck\) is part of the signed message, so the holder cannot change the PRF key between showings without
forging a new signature. The post-sign NIZK binds the public tag to the signed \(\veck\) while keeping \(\veck\) hidden.

\subsection{ARC security sketch}
\paragraph{Correctness.}
If issuance completes and the holder uses a fresh nonce, then the post-sign statement is satisfiable by construction,
by NIZK-ARK correctness (Definition~\ref{def:nizk-ark}) the verifier accepts.
\paragraph{Unforgeability and policy soundness.}
By knowledge soundness (Definition~\ref{def:nizk-ark}), any accepting presentation yields a witness
\((u,\mu,\veck,r_0,r_1)\) satisfying both the signature/hash relation and the PRF-tag relation.
Under one-more unforgeability of the blind signature (Section~\ref{sec:blind-signature}), an adversary cannot produce
accepting showings for credentials it did not obtain via issuance. Policy soundness is enforced by verifier state:
since \(\prftag=\PRF(\veck,\nonce)\) is uniquely determined by \((\veck,\nonce)\), two accepting
showings with the same nonce necessarily repeat the same tag, so a verifier that records spent pairs can reject replays.

\paragraph{Unlinkability.}
Presentation unlinkability is formalized in Definition~\ref{def:pres-unlink}. For distinct fresh nonces, PRF security
(Definition~\ref{def:prf}) makes tags pseudorandom, while zero knowledge of the post-sign proof (Definition~\ref{def:nizk-ark})
hides the credential secret. This yields indistinguishability of presentation transcripts, except for the intended
nonce-reuse caveat in Definition~\ref{def:pres-unlink}.

\paragraph{Attribute privacy and selective disclosure.}
The message part \(\mu\) can encode attributes; the post-sign NIZK can be extended to prove predicates about \(\mu\) and
to selectively reveal components, without changing the issuance protocol. We focus on the core rate-limiting
construction and defer predicate engineering to the SmallWood programming model in
Section~\ref{sec:smallwood-model}.
